%
% Copyright (C) 1995, 2000, 2013 American Mathematical Society.
% Copyright (C) 2016-2025 LaTeX Project and American Mathematical Society.
%
% This work may be distributed and/or modified under the
% conditions of the LaTeX Project Public License, either version 1.3c
% of this license or (at your option) any later version.
% The latest version of this license is in
%   https://www.latex-project.org/lppl.txt
% and version 1.3c or later is part of all distributions of LaTeX
% version 2005/12/01 or later.
%
% This work has the LPPL maintenance status `maintained'.
%
% The Current Maintainer of this work is the LaTeX Project.
%
\documentclass[leqno,titlepage,openany]{amsldoc}[1999/12/13]

\def\MaintainedByLaTeXTeam#1{%
\begin{center}%
\fbox{\fbox{\begin{tabular}{@{}l@{}}%
This file is maintained by the \LaTeX{} Project team.\\%
Bug reports can be opened (category \texttt{#1}) at\\%
\url{https://latex-project.org/bugs/}.\end{tabular}}}\end{center}}

\ifx\UndEfiNed\url
  \ClassError{amsldoc}{%
    This version of amsldoc.tex must be processed\MessageBreak
    with a newer version of amsldoc.cls (2.02 or later).}{}
\fi

\renewcommand\bibname{References}

%    Definitions that should be moved to amsldoc.cls
%    when it is next updated
\makeatletter
\let\cleardouble@page\cleardoublepage
\AtBeginDocument{%
  \ifx\cleardouble@page\cleardoublepage
    \def\cleardoublepage{\clearpage{\pagestyle{empty}\cleardouble@page}}
  \fi
}

% #1 is empty or one of class, environment, option, package
% could use different font in each case
\makeatletter
\def\@category@index#1#2{%
    \autoindex{%
        #2\actualchar\string\index#1font{#2}%
        \ifx\@nil#1\@nil\else\space#1\fi
    }%
}

\let\indexpackagefont\textsf
\let\indexclassfont\textsf
\let\indexoptionfont\texttt
\let\indexenvironmentfont\texttt
\let\indexfont\texttt
\makeatother
%    For a package that won't be indexed, i.e. amsmath
\newcommand{\nipkg}{\textsf}
\def\tikz/{Ti\textit{k}Z}
%    End definitions for amsldoc.cls

\title{User's Guide for the \nipkg{amsmath} Package (Version~2.1)}
\author{American Mathematical Society, \LaTeX\ Project}
\date{1999-12-13\\(revised 2002-02-25, 2016-11-14, 2018-04-05, 2019-10-14, 2020-02-18)}
\makeatletter
\def\@thanks{\bigskip\MaintainedByLaTeXTeam{amslatex}}
\makeatother

% Very minor changes are not listed in the \date above:
%
% 2021-08-22: reported typo fixed mathools -> mathtools
% 2025-05-23: drop spurious option list


%    Use the amsmath package and amscd package in order to print
%    examples.
\usepackage{amsmath}
\usepackage{amscd}

\usepackage{imakeidx}
%    Need alternate coding for vertical bar to get it into index.
%    Based on \qcbang and \cnbang in amsldoc.cls, refined by David Carlisle.
%    Give it a sort field that will be at the end of the symbols;
%    must avoid " and | which have special meanings.
\DeclareRobustCommand{\qcvert}{\qc\|\index{{z1}@{\ntt\char'174}}}
\DeclareRobustCommand{\cnvert}{%
  \ncn{\|}\index{{z0}@{\ntt\bslchar\char'174}}}
%    Need alternate index sort field for bang to allow & to sort first
\DeclareRobustCommand{\altcnbang}{%
  \ncn{\!}\index{"!b@{\ntt\bslchar\qcbang}}}
\DeclareRobustCommand{\cnamp}{%
  \index{"!a@\texttt{\&}}}

\makeindex[intoc] % generate index data
\providecommand{\see}[2]{\textit{see} #1}

%    The amsldoc class includes a number of features useful for
%    documentation about TeX, including:
%
%    ---Commands \tex/, \amstex/, \latex/, ... for uniform treatment
%    of the various logos and easy handling of following spaces.
%
%    ---Commands for printing various common elements: \cn for command
%    names, \fn for file names (including font-file names), \env for
%    environments, \pkg and \cls for packages and classes, etc.

%    Many of the command names used here are rather long and will
%    contribute to poor linebreaking if we follow the \latex/ practice
%    of not hyphenating anything set in tt font; instead we selectively
%    allow some hyphenation.
\allowtthyphens % defined in amsldoc.cls

\hyphenation{ac-cent-ed-sym-bol add-to-counter add-to-length align-at
  aligned-at allow-dis-play-breaks ams-art ams-cd ams-la-tex amsl-doc
  ams-symb ams-tex ams-text ams-xtra bmatrix bold-sym-bol cen-ter-tags
  eqn-ar-ray idots-int int-lim-its latex med-space neg-med-space
  neg-thick-space neg-thin-space no-int-lim-its no-name-lim-its
  over-left-arrow over-left-right-arrow over-right-arrow pmatrix
  qed-sym-bol set-length side-set small-er tbinom the-equa-tion
  thick-space thin-space un-der-left-arrow un-der-left-right-arrow
  un-der-right-arrow use-pack-age var-inj-lim var-proj-lim vmatrix
  xalign-at xx-align-at}

%    Prepare for illustrating the \vec example
\newcommand{\vect}[1]{\mathbf{#1}}

\newcommand{\booktitle}[1]{\textit{#1}}
\newcommand{\journalname}[1]{\textit{#1}}
\newcommand{\seriesname}[1]{\textit{#1}}

%    Command to insert and index a particular phrase. Doesn't work for
%    certain kinds of special characters in the argument.
\newcommand{\ii}[1]{#1\index{#1}}

\newcommand{\vstrut}[1]{\vrule width0pt height#1\relax}

%    An environment for presenting comprehensive address information:
\newenvironment{infoaddress}{%
  \par\topsep\medskipamount
  \trivlist\centering
  \item[]%
  \begin{minipage}{.7\columnwidth}%
  \raggedright
}{%
  \end{minipage}%
  \endtrivlist
}

\newenvironment{eqxample}{%
  \par\addvspace\medskipamount
  \noindent\begin{minipage}{.5\columnwidth}%
  \def\producing{\end{minipage}\begin{minipage}{.5\columnwidth}%
    \hbox\bgroup\kern-.2pt\vrule width.2pt%
      \vbox\bgroup\parindent0pt\relax
%    The 3pt is to cancel the -\lineskip from \displ@y
    \abovedisplayskip3pt \abovedisplayshortskip\abovedisplayskip
    \belowdisplayskip0pt \belowdisplayshortskip\belowdisplayskip
    \noindent}
}{%
  \par
%    Ensure that a lonely \[\] structure doesn't take up width less than
%    \hsize.
  \hrule height0pt width\hsize
  \egroup\vrule width.2pt\kern-.2pt\egroup
  \end{minipage}%
  \par\addvspace\medskipamount
}

%    Errors in output shouldn't be \texttt, and they may also be longer
%    than one line, but they should be identified by the "errorbullet".
%    (They were formerly labeled as aubsections.)
\newenvironment{erroro}[1]{%
  \par\addvspace\medskipamount
  \noindent\hangindent\parindent
  \hbox to\parindent{\errorbullet\hfil}\ignorespaces
  #1\par\vspace{\smallskipamount}
}{%
  \par\addvspace\medskipamount
}

%    The chapters are so short, perhaps we shouldn't call them by the
%    name `Chapter'. We make \chaptername read an argument in order to
%    remove a following \space or "{} " (both possibilities are present
%    in book.cls).

\renewcommand{\chaptername}[1]{}
\newcommand{\chapnum}[1]{\mdash #1\mdash }
\makeatletter
\def\@makechapterhead#1{%
  \vspace{1.5\baselineskip}%
  {\parindent \z@ \raggedright \reset@font
    \ifnum \c@secnumdepth >\m@ne
      \large\bfseries \chapnum{\thechapter}%
      \par\nobreak
      \vskip.5\baselineskip\relax
    \fi
    #1\par\nobreak
    \vskip\baselineskip
  }}
\makeatother

%    A command for ragged-right parbox in a tabular.
\newcommand{\rp}{\let\PBS\\\raggedright\let\\\PBS}

%    Non-indexed file name
\newcommand{\nfn}[1]{\texttt{#1}}

%    For the examples in the math spacing table.
%%\newcommand{\lspx}{\mbox{\rule{5pt}{.6pt}\rule{.6pt}{6pt}}}
%%\newcommand{\rspx}{\mbox{\rule[-1pt]{.6pt}{7pt}%
%%  \rule[-1pt]{5pt}{.6pt}}}
\newcommand{\lspx}{\mathord{\Rightarrow\mkern-1mu}}
\newcommand{\rspx}{\mathord{\mkern-1mu\Leftarrow}}
\newcommand{\spx}[1]{$\lspx #1\rspx$}

%    For a list of characters representing document input.
\newcommand{\clist}[1]{%
  \mbox{\ntt\spaceskip.2em plus.1em \xspaceskip\spaceskip#1}}

%    Fix weird \latex/ definition of rightmark.
\makeatletter
\def\rightmark{\expandafter\@rightmark\botmark{}{}}
%    Also turn off section marks.
\let\sectionmark\@gobble
\renewcommand{\chaptermark}[1]{%
  \uppercase{\markboth{\rhcn#1}{\rhcn#1}}}
\newcommand{\rhcn}{\thechapter. }
\makeatother

%    Include down to \section but not \subsection, in toc:
\setcounter{tocdepth}{1}

%    Subheadings for the bibliography
\newif\iffirstbibsubhead \firstbibsubheadtrue
\newcommand{\bibsubhead}[1]{%
  \iffirstbibsubhead \firstbibsubheadfalse
  \else \addvspace{\medskipamount}
  \fi
  \item[]\hspace*{-\leftmargin}\textbf{#1}\par
  \vspace{\smallskipamount}
}

\DeclareMathOperator{\ix}{ix}
\DeclareMathOperator{\nul}{nul}
\DeclareMathOperator{\End}{End}
\DeclareMathOperator{\xxx}{xxx}

\usepackage{hyperref}
\hypersetup{pdfborder=0 0 0}
\makeatletter
\let\oldcs\cs
\def\cs#1{\texorpdfstring{\oldcs{#1}}{\@backslashchar\@backslashchar#1}}
\let\cn\cs
\makeatother
\begin{document}
\section{Alignment building blocks}

Like \env{equation}, the multi-equation environments \env{gather},
\env{align}, and \env{alignat} are designed to produce a structure
whose width is the full line width. This means, for example, that one
cannot readily add parentheses around the entire structure. But
variants\index{ed env@\texttt{-ed} environments|(}
\env{gathered}, \env{aligned}, and \env{alignedat} are provided whose
total width is the actual width of the contents; thus they can be used
as a component in a containing expression. E.g.,
\begin{equation*}
\left.\begin{aligned}
  B'&=-\partial\times E,\\
  E'&=\partial\times B - 4\pi j,
\end{aligned}
\right\}
\qquad \text{Maxwell's equations}
\end{equation*}
\begin{verbatim}
\begin{equation*}
\left.\begin{aligned}
  B'&=-\partial\times E,\\
  E'&=\partial\times B - 4\pi j,
\end{aligned}
\right\}
\qquad \text{Maxwell's equations}
\end{equation*}
\end{verbatim}
Like the \env{array} environment, these \texttt{-ed} variants also take
an optional\index{options!positioning of \texttt{-ed} environments}
|[t]|\index{t option@\texttt{t} (top) option},
|[b]|\index{b option@\texttt{b} (bottom) option} or the default
|[c]|\index{c option@\texttt{c} (center) option}
argument to specify vertical
positioning.  For maximum interoperability, do not insert a
space\index{options!space before \texttt{[}}
or line break\index{line break} before the option.  See
also the note on page~\pageref{breaknote} regarding page breaking
within the \texttt{-ed} environments.

\qq{Cases} constructions like the following are common in
mathematics:
\begin{equation}\label{eq:C}
P_{r-j}=
  \begin{cases}
    0&  \text{if $r-j$ is odd},\\
    r!\,(-1)^{(r-j)/2}&  \text{if $r-j$ is even}.
  \end{cases}
\end{equation}
and in the \nipkg{amsmath} package there is a \env{cases} environment to
make them easy to write:
\begin{verbatim}
P_{r-j}=\begin{cases}
    0&  \text{if $r-j$ is odd},\\
    r!\,(-1)^{(r-j)/2}&  \text{if $r-j$ is even}.
  \end{cases}
\end{verbatim}
Notice the use of \cn{text} (cf.~\secref{text}) and the nested math
formulas. \env{cases} is set in \cn{textstyle}. If \cn{displaystyle}
is wanted, it must be requested explicitly; \pkg{mathtools} provides
a \env{dcases} environment for this purpose.

The |-ed|\index{ed env@\texttt{-ed} environments|)} and \env{cases}
environments must appear within an enclosing math environment, which can
be either in text, between |$...$|, or in any of the display environments.

%% ------------------------------------------------------------------ %%

\section{Adjusting tag placement}

Placing equation numbers can be a rather complex problem in multiline
displays. The environments of the \nipkg{amsmath} package try hard to
avoid overprinting an equation number on the equation contents, if
necessary moving the number down or up to a separate line. Difficulties
in accurately calculating the profile of an equation can occasionally
result in number movement that doesn't look right.
A \cn{raisetag} command is provided to adjust the vertical position of the
current equation number, if it has been shifted away from its normal
position. If it is in its normal position the command has no effect whatsoever.
In particular, this means that one can't use it to shift the tag of an \env{equation}
environment.

To move a particular number up by six points, write
|\raisetag{6pt}|. (At the end of a display, this also shifts up the
text following the display.)
This kind of adjustment is fine tuning like line
breaks and page breaks, and should therefore be left undone until your
document is nearly finalized, or you may end up redoing the fine tuning
several times to keep up with changing document contents.

%% ------------------------------------------------------------------ %%

\section{Vertical spacing and page breaks in multiline displays}

You can use the \cn{\\}|[|\<dimension>|]|%
\index{options!extra vertical space after \cn{\\}}
command to get extra vertical
space between lines in all the \nipkg{amsmath} displayed equation
environments, as is usual in \latex/. Do not type a space between the%
\index{options!space before \texttt{[}}
\cn{\\} and the following |[|; \emph{only} for display environments
defined by \nipkg{amsmath} the space is interpreted to mean that the
bracketed material is part of the document content.

When the \nipkg{amsmath} package is
in use \ii{page breaks} between equation lines are normally disallowed;
the philosophy is that page breaks in such material should receive
individual attention from the author. To get an individual page break
inside a particular displayed equation, a \cn{displaybreak} command is
provided. \cn{displaybreak} is best placed immediately before the
\cn{\\} where it is to take effect.  Like \latex/'s \cn{pagebreak},
\cn{displaybreak} takes an optional argument%
\index{options!behavior of particular options} between 0 and 4 denoting
the desirability of the pagebreak. |\displaybreak[0]| means \qq{it is
  permissible to break here} without encouraging a break;
\cn{displaybreak} with no optional argument is the same as
|\displaybreak[4]| and forces a break.

If you prefer a strategy of letting page breaks fall where they may,
even in the middle of a multiline equation, then you might put
\cn{allowdisplaybreaks}\texttt{[1]} in the preamble of your document. An
optional argument\index{options!behavior of particular options}
1\ndash 4 can be used for finer control: |[1]| means
allow page breaks, but avoid them as much as possible; values of 2,3,4
mean increasing permissiveness. When display breaks are enabled with
\cn{allowdisplaybreaks}, the \cn{\\*} command can be used to prohibit a
pagebreak after a given line, as usual.

\begin{bfseries}
Note:\label{breaknote} Certain equation environments wrap their contents
in an unbreakable box, with the consequence that neither \cn{displaybreak}
nor \cn{allowdisplaybreaks} will have any effect on them. These include
\env{split}, \env{aligned}, \env{gathered}, and \env{alignedat}.
\end{bfseries}

%% ------------------------------------------------------------------ %%

\section{Interrupting a display}

The command \cn{intertext} is used for a short interjection of one or
two lines of text\index{text fragments inside math} in the middle of a
multiple-line display structure (see also the \cn{text} command in
\secref{text}). Its salient feature is preservation of the alignment,
which would not happen if you simply ended the display and then started
it up again afterwards. \cn{intertext} may only appear right after a
\cn{\\} or \cn{\\*} command. Notice the position of the word \qq{and} in
this example.
\begin{align}
  A_1&=N_0(\lambda;\Omega')-\phi(\lambda;\Omega'),\\
  A_2&=\phi(\lambda;\Omega')-\phi(\lambda;\Omega),\\
\intertext{and}
  A_3&=\mathcal{N}(\lambda;\omega).
\end{align}
\begin{verbatim}
\begin{align}
  A_1&=N_0(\lambda;\Omega')-\phi(\lambda;\Omega'),\\
  A_2&=\phi(\lambda;\Omega')-\phi(\lambda;\Omega),\\
\intertext{and}
  A_3&=\mathcal{N}(\lambda;\omega).
\end{align}
\end{verbatim}

The \pkg{mathtools} package provides a command \cn{shortintertext} that
is intended for use when the interjected text is only a few words; it uses
less vertical space than \cn{intertext}. This
is most effective when equation numbers are on the right.%
\index{equationn@equation numbers!left or right placement}

%% ------------------------------------------------------------------ %%

\section{Equation numbering}

\subsection{Numbering hierarchy}
In \latex/ if you wanted to have equations numbered within
sections\mdash that is, have
equation numbers (1.1), (1.2), \dots, (2.1), (2.2),
\dots, in sections 1, 2, and so forth\mdash you could redefine
\cn{theequation} as suggested in the \latex/ manual \cite[\S6.3, \S
C.8.4]{lamport}:
\begin{verbatim}
\renewcommand{\theequation}{\thesection.\arabic{equation}}
\end{verbatim}

This works pretty well, except that the equation counter won't be reset
to zero at the beginning of a new section or chapter, unless you do it
yourself using \cn{setcounter}. To make this a little more convenient,
the \nipkg{amsmath} package provides a
command\index{equationn@equation numbers!hierarchy} \cn{numberwithin}.
To have equation numbering tied to section numbering, with automatic
reset of the equation counter, write
\begin{verbatim}
\numberwithin{equation}{section}
\end{verbatim}
As its name implies, the \cn{numberwithin} command can be applied to
any counter, not just the \texttt{equation} counter.

\subsection{Cross references to equation numbers}

To make cross-references to equations easier, an \cn{eqref}
command\index{equationn@equation numbers!cross-references} is provided. This
automatically supplies the parentheses around the equation number. I.e.,
if \verb'\ref{abc}' produces 3.2 then \verb'\eqref{abc}' produces
(3.2). The parentheses around an \cn{eqref} equation number will be
set in upright type regardless of type style of the context.

\subsection{Subordinate numbering sequences}

The \nipkg{amsmath} package provides also a wrapper environment,
\env{subequations}\index{equationn@equation numbers!subordinate numbering}, to make it
easy to number equations in a particular \cn{align} or similar group with
a subordinate numbering scheme. For example
\begin{verbatim}
\begin{subequations}
...
\end{subequations}
\end{verbatim}
causes all numbered equations within that part of the document to be
numbered (4.9a) (4.9b) (4.9c) \dots, if the preceding numbered
equation was (4.8). A \cn{label} command immediately after
\verb/\begin{subequations}/ will produce a \cn{ref} of the parent
number 4.9, not 4.9a. The counters used by the subequations
environment are \verb/parentequation/ and \verb/equation/ and
\cn{addtocounter}, \cn{setcounter}, \cn{value}, etc., can be applied
as usual to those counter names. To get anything other than lowercase
letters for the subordinate numbers, use standard \latex/ methods for
changing numbering style \cite[\S6.3, \S C.8.4]{lamport}. For example,
redefining \cn{theequation} as follows will produce roman numerals.
\begin{verbatim}
\begin{subequations}
\renewcommand{\theequation}{\theparentequation \roman{equation}}
...
\end{verbatim}

\subsection{Numbering style}

The default equation number is set in \cn{normalfont}.  This means that
in bold section headings, bold is suppressed; a workaround for this is
to use \verb+(\ref{+\ldots\verb+})+ rather than \verb+\eqref{+\ldots\verb+}+.

If smaller type is specified for a numbered display, the size of the
equation number will also be small.  The default size can be ensured
throughout a document by applying this patch in the preamble:
\begin{verbatim}
\makeatletter
\renewcommand{\maketag@@@}[1]{\hbox{\m@th\normalsize\normalfont#1}}%
\makeatother
\end{verbatim}
(This modification may be included in a future version of
\nipkg{amsmath}.)


%%%%%%%%%%%%%%%%%%%%%%%%%%%%%%%%%%%%%%%%%%%%%%%%%%%%%%%%%%%%%%%%%%%%%%%%

\end{document}
